\section{Compact Riemann Surfaces}
\label{sec:compact-riemann-surfaces}

Throughout, \(X,Y,Z\) denote compact connected Riemann surfaces, and maps
between them are holomorphic unless otherwise specified. Compactness implies
that every holomorphic function \(X\to\C\) is constant. Holomorphic maps
between compact Riemann surfaces are either constant or finite branched
coverings.

\begin{classicalinput}[Meromorphic functions and divisors]
\label{input:divisors}
There is a divisor group \(\Div(X)\), a subgroup of principal divisors
\(\Prin(X)\), a degree map \(\deg:\Div(X)\to\Z\), and the Riemann-Roch spaces
\[
  L(D)=\{ f\text{ meromorphic on }X : (f)+D\ge 0\}.
\]
Principal divisors have degree zero, and a nonconstant meromorphic function
defines a holomorphic map \(X\to\CPone\).
\end{classicalinput}

\begin{proofsketch}
The local theory of holomorphic functions gives zeros and poles with finite
orders. Summing orders defines divisors. Compactness gives finiteness of the
support of a divisor associated to a meromorphic function. The meromorphic
function is interpreted as a holomorphic map to \(\CPone\) by using the value
\(\infty\) at poles.
\end{proofsketch}

\begin{formalizationnote}
The repository currently has substantial chart and path-integration
infrastructure, but not a full divisor or meromorphic-function theory on
compact Riemann surfaces. This input is the source for later theorem names
around Riemann-Roch, Abel's theorem, and degree theory.
\end{formalizationnote}

